\section{System Design}


\subsection{Architecture Overview}

Verus comprises three layers:

\begin{enumerate}
\item \textbf{Ingestion layer}: Receives events, validates, and routes to processing.
\item \textbf{Processing layer}: Maintains real-time aggregations and writes to storage.
\item \textbf{Query layer}: Routes queries to optimal backend and merges results.
\end{enumerate}

\begin{figure}[h]
\centering
\fbox{\parbox{0.9\linewidth}{
\centering
\textit{[Architecture diagram]}\\[0.5em]
Events $\rightarrow$ Kafka $\rightarrow$ Stream Processor $\rightarrow$ Redis (aggregates)\\
\hspace{4em} $\searrow$ Columnar Store (raw events)\\[0.5em]
Query Router $\rightarrow$ Redis / Columnar / Both
}}
\caption{Verus architecture overview}
\label{fig:arch}
\end{figure}

\subsection{Ingestion Layer}

Events arrive via HTTP endpoints or Kafka consumers:

\begin{lstlisting}[language=Python,caption=Event ingestion]
class EventIngester:
    def __init__(self, kafka_producer, validator):
        self.producer = kafka_producer
        self.validator = validator

    async def ingest(self, event: CommerceEvent) -> None:
        # Validate schema and business rules
        self.validator.validate(event)

        # Enrich with server-side data
        event.server_timestamp = datetime.utcnow()
        event.geo = geoip.lookup(event.ip_address)

        # Partition by store for ordering guarantees
        partition_key = event.store_id

        await self.producer.send(
            topic="commerce-events",
            key=partition_key,
            value=event.to_json()
        )
\end{lstlisting}

Events partition by store ID, ensuring ordering within a store while enabling parallel processing across stores.

\subsection{Stream Processing Layer}

The stream processor maintains real-time aggregations:

\subsubsection{Aggregation Windows}

We maintain aggregations at multiple time granularities:

\begin{itemize}
\item \textbf{1-minute windows}: Operational dashboards.
\item \textbf{1-hour windows}: Intraday trends.
\item \textbf{1-day windows}: Daily summaries.
\end{itemize}

Windows use event time with allowed lateness of 5 minutes. Late events update existing windows; very late events are logged but excluded.

\subsubsection{Materialized Views}

Common query patterns are pre-computed as materialized views:

\begin{lstlisting}[language=SQL,caption=Materialized view definition]
CREATE MATERIALIZED VIEW revenue_by_source AS
SELECT
    store_id,
    TUMBLE(event_time, INTERVAL '1' MINUTE) AS window,
    utm_source,
    COUNT(*) AS order_count,
    SUM(revenue) AS total_revenue,
    COUNT(DISTINCT user_id) AS unique_customers
FROM events
WHERE event_type = 'purchase'
GROUP BY store_id, window, utm_source;
\end{lstlisting}

The stream processor maintains these views incrementally:

\begin{lstlisting}[language=Python,caption=Incremental aggregation]
class MaterializedViewProcessor:
    def __init__(self, view_def, state_store):
        self.view_def = view_def
        self.state = state_store

    def process(self, event):
        if not self.view_def.filter(event):
            return

        # Compute aggregation key
        key = self.view_def.group_key(event)

        # Update state
        current = self.state.get(key) or self.view_def.initial_state()
        updated = self.view_def.aggregate(current, event)
        self.state.put(key, updated)

        # Emit to downstream
        self.emit(key, updated)
\end{lstlisting}

\subsubsection{State Management}

Aggregation state resides in Redis for fast access:

\begin{lstlisting}[language=Python,caption=Redis state store]
class RedisStateStore:
    def __init__(self, redis_client, ttl_seconds):
        self.redis = redis_client
        self.ttl = ttl_seconds

    def get(self, key: str) -> Optional[Dict]:
        data = self.redis.get(key)
        return json.loads(data) if data else None

    def put(self, key: str, value: Dict) -> None:
        self.redis.setex(key, self.ttl, json.dumps(value))

    def increment(self, key: str, field: str, delta: float) -> float:
        return self.redis.hincrbyfloat(key, field, delta)
\end{lstlisting}

State keys encode view name, dimensions, and time window:

\begin{verbatim}
verus:revenue_by_source:store_123:google:2015-06-15T14:30
\end{verbatim}

\subsection{Columnar Storage Layer}

Raw events persist to columnar storage for ad-hoc queries:

\subsubsection{Storage Format}

We use Apache Parquet with the following optimizations:

\begin{itemize}
\item \textbf{Partitioning}: By store ID and date for partition pruning.
\item \textbf{Sorting}: By timestamp within partitions for range queries.
\item \textbf{Compression}: Snappy for balance of speed and size.
\item \textbf{Statistics}: Min/max per column for predicate pushdown.
\end{itemize}

\begin{lstlisting}[language=Python,caption=Parquet writer]
class ParquetWriter:
    def __init__(self, base_path, partition_cols):
        self.base_path = base_path
        self.partition_cols = partition_cols
        self.buffers = {}  # partition -> event buffer

    def write(self, event):
        partition = self.compute_partition(event)
        self.buffers.setdefault(partition, []).append(event)

        if len(self.buffers[partition]) >= BATCH_SIZE:
            self.flush(partition)

    def flush(self, partition):
        events = self.buffers.pop(partition)
        path = f"{self.base_path}/{partition}/data_{uuid4()}.parquet"

        table = pa.Table.from_pylist(events, schema=EVENT_SCHEMA)
        pq.write_table(table, path, compression='snappy')
\end{lstlisting}

\subsubsection{Query Engine}

Ad-hoc queries execute on columnar storage via a distributed query engine (we use Presto):

\begin{lstlisting}[language=SQL,caption=Ad-hoc query example]
SELECT
    product_category,
    COUNT(*) AS views,
    COUNT(DISTINCT session_id) AS sessions,
    SUM(CASE WHEN event_type = 'purchase' THEN 1 ELSE 0 END) AS purchases
FROM events
WHERE store_id = 'store_123'
  AND event_date = '2015-06-15'
  AND device_type = 'mobile'
GROUP BY product_category
ORDER BY views DESC
LIMIT 20;
\end{lstlisting}

\subsection{Query Routing Layer}

The query router examines incoming queries and routes to the optimal backend:

\begin{lstlisting}[language=Python,caption=Query routing]
class QueryRouter:
    def __init__(self, redis_client, presto_client, view_catalog):
        self.redis = redis_client
        self.presto = presto_client
        self.views = view_catalog

    def execute(self, query: str) -> ResultSet:
        parsed = parse_sql(query)

        # Check if query matches a materialized view
        view_match = self.views.match(parsed)

        if view_match and self.is_recent_enough(parsed, view_match):
            # Serve from Redis
            return self.query_redis(view_match, parsed)
        elif view_match:
            # Merge Redis (recent) + Presto (historical)
            return self.merge_results(
                self.query_redis(view_match, parsed),
                self.query_presto(parsed)
            )
        else:
            # Full scan on Presto
            return self.query_presto(parsed)

    def is_recent_enough(self, parsed, view):
        # Check if query time range fits in Redis TTL
        query_start = parsed.get_time_range()[0]
        redis_start = datetime.utcnow() - timedelta(seconds=view.ttl)
        return query_start >= redis_start
\end{lstlisting}

